\section{Discussion}
\label{sec:discussion}

\subsection{What the results imply for practice}
\label{sec:discussion-practice}

The diagnostics in Section~\ref{sec:decomposition} place the ladder results in
context, but do not identify a causal mechanism.

\paragraph{Why standardisation captures the observed gain.} The ladder and its one-sided bounds
show that per-dimension standardisation captures the observed improvement: no
more complex rung is shown to improve target macro-F1 in either direction.
The evaluated CORAL and affine MK-MMD maps do not yield a further gain. The
diagnostics report raw marginal and class-conditional MMD$^2$ in one
source-defined reference geometry. Their class-specific null effect sizes aid
finite-sample interpretation but are not divided by the marginal effect size.
The fixed-bandwidth z-score control saturates the RBF kernel, so it is reported
as a measurement limitation rather than a closeness result.

\paragraph{Why lower marginal discrepancy does not help further here.}
The MMD rungs can reduce their fitted marginal objective without improving the
task metric. The class-conditional diagnostic provides a descriptive account of
that mismatch, but it is not a direct estimate of $P(y \mid x)$.

This analysis is limited to these rungs and this corpus pair. Covariate-shift
correction is justified by the assumption that $P(y \mid x)$ is stable across
domains. Our diagnostics do not directly test that assumption. Observed
label-prior differences are small: the prior KL is $0.02235$ and $0.02541$
nats, and the standard corrections examined here \emph{hurt} in 238 of 240
cases, by $0.013$ to $0.24$ macro-F1, while the same estimators recover a
planted shift. This does not establish the absence of label shift, because these
corrections rely on class-conditional assumptions that the diagnostic does not
verify.

\paragraph{What a correct selection protocol would have to measure.} The
forward selection result establishes a specific failure: source-side validation
does not reliably separate unaligned from standardised features in that
direction. It does not establish the cause. The class-conditional diagnostic
motivates a target-side signal sensitive to variation beyond the marginal
objective, but it is not itself such a signal and it requires target labels.
We do not have an unsupervised alternative to offer. What any candidate signal
must do is separate configurations that source-side validation cannot, and
Section~\ref{sec:results-selection} shows how wide that gap is.

\subsection{Selection is where this actually bites}
\label{sec:discussion-selection}

Table~\ref{tab:headline} gives the practitioner-facing validated scores and the
oracle gaps. \textbf{The forward oracle gap is larger than the observed
alignment gain.} It describes the room between source-side selection and the
target-set maximum; it does not isolate a causal effect of source-side
selection.

The forward selection surface is specific rather than universal: source-side
validation is nearly blind to the alignment distinction in that direction and
selects an unaligned, below-chance configuration once. A practitioner following
the protocol uses no target labels and has no source-side indication of that
failure. By contrast, source-side validation is informative on the layer-depth
axis in a separate analysis; the present result is limited to the alignment
decision.

This is also the argument for reporting baselines beside every performance
number. Without the chance baseline, a below-baseline result reads only as a
weak score.

\subsection{Testing the conditional-shift explanation directly}
\label{sec:discussion-testing}

Our class-conditional diagnostic is not a direct test of conditional shift. It
compares $P(x \mid y)$ across corpora, whereas the covariate-shift assumption
concerns $P(y \mid x)$. The falsifiable prediction we could test; that
label-shift correction should not help under near-zero prior KL; passed,
and passed sharply. That rules out neither a conditional-shift mechanism nor
other explanations for the plateau, so we state what a direct test would take.

It would need paired data: the same speakers, or the same sentences, recorded
under both corpora's conditions, so that $P(y \mid x)$ can be estimated
separately in each domain without the corpus difference and the speaker
difference being confounded. Failing that, per-utterance perceptual ratings from
multiple annotators would let the conditional be compared where the label is a
distribution rather than a point. Neither exists for this pair, so the
diagnostic here motivates a direct test rather than supplying one.

\subsection{Limitations}
\label{sec:limitations}

\textbf{Two corpora, one language, two directions.} Every result is a statement
about RAVDESS and CREMA-D in English. We have not shown that the observed
pattern, marginal discrepancy declining without a further transfer gain,
appears elsewhere. A corpus pair with a genuinely large prior difference
would test the label-shift correction result much harder than this one does.

\textbf{Both corpora are acted.} Neither contains spontaneous speech. Acted
emotion is more prototypical and more clearly separated than spontaneous
emotion, so the absolute scores here are likely optimistic relative to
spontaneous data. Whether the observed pattern transfers is genuinely open:
acted and spontaneous corpora may differ in ways not represented by this pair.

\textbf{Modest absolute performance.} Validated macro-F1 of $0.3156$ and
$0.4994$ clears chance comfortably but falls well short of in-domain SER. We
report a controlled analysis, not a system, and nothing here should be read as a
recommended pipeline.

\textbf{The reverse direction is less precisely measured.} RAVDESS as target
supplies $624$ test utterances against CREMA-D's $3677$ to $3690$, so reverse
intervals are wider throughout; for the among-rung bound, about $1.8$ times.
Both directions are matched in source-training size only. Their
source-validation and target-adaptation sizes still differ, so source-training
size cannot explain the directional asymmetry but other data-quantity effects
remain possible; the asymmetry in \emph{precision} is directly explained by the
different target-test sizes.

\textbf{The Transformer arm is reduced-seed.} Two seeds, the primary direction
only, one inner-grid setting per rung. It is never pooled with the five-seed
families and its intervals are correspondingly wider.

\textbf{The MK-MMD rungs frequently fall back.} Optimisation reverts to its
warm start in $64.9\%$ of \texttt{mkmmd\_diag} runs and $55.2\%$ of
\texttt{mkmmd\_full} runs, so in the majority of cells the fitted map \emph{is}
the CORAL or diagonal warm start. This qualifies the negative result about the
evaluated affine MK-MMD maps: this optimiser, at this budget, does not beat
standardisation. A better-optimised MK-MMD map might do better; these runs do
not decide that question.
